\begin{definition}
    \textbf{Limite finito di una successione}

    Si dice che il limite di una successione (il valore a cui tende per
    $n\rightarrow+\infty$) è finito se e solo se la successione è convergente.
    
    Se il limite di una successione è finito, avremo che il limite $l\in\mathbb{R}$
    di quella successione coincide con il valore di convergenza.

    Dunque, in simboli,
    \[
        \underset{n\rightarrow+\infty}{\lim} x_n = l \Leftrightarrow
        \exists l \in\mathbb{R} \mid \forall\epsilon>0, \exists n_\epsilon\in\mathbb{N}\mid
        \forall n \geq n_\epsilon, \left|x_n-l\right|<\epsilon
    \]
    oppure, equivalentemente,
    \[
        \underset{n\rightarrow+\infty}{\lim} x_n = l \Leftrightarrow
        \exists l \in\mathbb{R} \mid \forall\epsilon>0, \exists n_\epsilon\in\mathbb{N}\mid
        \forall n \geq n_\epsilon, x_n\in\left(l-\epsilon, l+\epsilon\right).
    \]
\end{definition}

\begin{observation}
    Una serie convergente ammette esattamente un solo valore di convergenza, dunque
    ogni serie con un limite ammette un solo valore di convergenza, che corrisponderà
    con il limite di quella successione. Il limite finito di una successione è dunque
    unico poiché è unico il valore di convergenza della serie.
\end{observation}

Analogamente definiamo i seguenti limiti particolari

\begin{definition}
    \textbf{Limite destro finito di una successione}

    Sia $\underset{n\rightarrow+\infty}{\lim} x_n = l$, allora si può specificare
    che $x_n$ converge a $l$ da destra e scrivendo che
    \[
        \underset{n\rightarrow+\infty}{\lim} x_n = l^+ \Leftrightarrow
        \exists l \in\mathbb{R} \mid \forall\epsilon>0, \exists n_\epsilon\in\mathbb{N}\mid
        \forall n \geq n_\epsilon, x_n\in\left[l, l+\epsilon\right).
    \]
\end{definition}

\begin{definition}
    \textbf{Limite sinsitro finito di una successione}

    Sia $\underset{n\rightarrow+\infty}{\lim} x_n = l$, allora si può specificare
    che $x_n$ converge a $l$ da sinistra e scrivendo che
    \[
        \underset{n\rightarrow+\infty}{\lim} x_n = l^- \Leftrightarrow
        \exists l \in\mathbb{R} \mid \forall\epsilon>0, \exists n_\epsilon\in\mathbb{N}\mid
        \forall n \geq n_\epsilon, x_n\in\left(l-\epsilon, l\right].
    \]
\end{definition}

\begin{observation}
    Il fatto che la successioni converga da destra o da sinsitra a un certo valore
    non cambia il valore effettivo del limite che resta il valore di convergenza.
    Aggiungere un segno ad apice è facoltativo e serve semplicemente a evidenziare che
    la successione converge da destra o da sinsitra. Il valore numerico del limite
    rimane sempre unico, difatti dire che una successione converge a un valore da una
    particolare direzione implica necessariamente che la successione converga a
    quel valore specifico, in simboli

    \[
        \begin{cases}
            \underset{n\rightarrow+\infty}{\lim} x_n = l^-\implies\underset{n\rightarrow+\infty}{\lim} x_n = l \\
            \underset{n\rightarrow+\infty}{\lim} x_n = l^+\implies\underset{n\rightarrow+\infty}{\lim} x_n = l
        \end{cases}.
    \]
\end{observation}

\begin{example}\leavevmode

    \begin{enumerate}[label=(\roman*)]
        \item \[
            \frac{1}{n} \rightarrow 0^+ \left(\underset{n\rightarrow+\infty}{\lim}\frac{1}{n}=0^+\right)
        \]
        \item \[
            1-\frac{1}{n} \rightarrow 1^-\left(\underset{n\rightarrow+\infty}{\lim}1-\frac{1}{n}=1^-\right)
        \]
        \item \[
            1-\frac{\left(-1\right)^n}{n} \rightarrow 1
            \left(\underset{n\rightarrow+\infty}{\lim}\left[1-\frac{\left(-1\right)^n}{n}\right]=
            1\right)
        \]
    \end{enumerate}
\end{example}

\begin{observation}
    Dalla definizone segue direttamente che
    \[
        \underset{n\rightarrow+\infty}{\lim} x_n = l \Leftrightarrow
        x_n \rightarrow l \Leftrightarrow
        x_n-l \rightarrow 0 \Leftrightarrow
        \left|x_n-l\right| \rightarrow 0.
    \]
    In particolare con $l=0$
    \[
        \underset{n\rightarrow+\infty}{\lim} x_n = 0 \Leftrightarrow
        \left|x_n - 0\right| \rightarrow 0
        \left|x_n\right| \rightarrow 0.
    \]
\end{observation}

\begin{corollary}\label{cor:con-val-ass}
    Se una serie converge a un valore, allora il suo valore assoluto converge al valore
    assoluto del valore di convergenza iniziale (NON vale necessariamente l'inversa).
    In simboli
    \[
        x_n\rightarrow l\implies \left|x_n\right|\rightarrow\left|l\right|.
    \]
\end{corollary}

\begin{demonstration}
    \textbf{di} \corref{cor:con-val-ass}

    [dimostrazione]
\end{demonstration}